\chapter{Architecture}
\label{ch:architecture}

Figure~\ref{fig:architecture} shows the agents, the stores they read and write, the human checkpoints and the PM tool. It is redrawn from the full Mermaid diagram in \code{docs/design/architecture.mmd}.

\begin{figure}[htbp]
  \centering
  \resizebox{\textwidth}{!}{% Architecture of the Harbourline AI PM system (redrawn from docs/design/architecture.mmd)
\begin{tikzpicture}[
    font=\footnotesize,
    >={Stealth[length=2.2mm]},
    box/.style={rectangle, rounded corners=2pt, align=center, text width=3.05cm,
                minimum height=1.0cm, inner sep=3pt},
    agent/.style={box, draw=pmNavy, very thick, fill=white},
    src/.style={box, draw=pmGrey, fill=pmLight},
    store/.style={box, draw=pmBlue, fill=pmBlue!10},
    llm/.style={box, draw=pmGrey, dashed, fill=white},
    ext/.style={box, draw=pmTeal, thick, fill=pmTeal!8},
    human/.style={box, draw=hcLine, dashed, thick, fill=white},
    hcin/.style={align=center, text width=2.0cm, inner xsep=0pt, inner ysep=3pt,
                 minimum height=0.95cm, font=\footnotesize},
    flow/.style={->, thick, pmNavy!80!black},
    hand/.style={->, line width=2.2pt, pmOrange},
    dot/.style={->, densely dotted, thick, pmGrey!80!black},
    resetline/.style={->, dashed, thick, hcLine},
    lab/.style={font=\scriptsize, align=center, inner sep=1pt, fill=white},
]
% hexagonal human checkpoint: \hc{name}{position}{text}
\newcommand{\hc}[3]{%
  \node[hcin] (#1) at (#2) {#3};
  \begin{scope}[on background layer]
    \draw[fill=hcFill, draw=hcLine, thick]
      ($(#1.west)+(-0.32,0)$) -- (#1.north west) -- (#1.north east) --
      ($(#1.east)+(0.32,0)$) -- (#1.south east) -- (#1.south west) -- cycle;
  \end{scope}
  \coordinate (#1-w) at ($(#1.west)+(-0.32,0)$);
  \coordinate (#1-e) at ($(#1.east)+(0.32,0)$);
}

% ---------- sources and fact store ----------
\node[src]   (PACK)  at (0,0)      {\textbf{Data pack 00--07}\\brief + sponsor answers};
\node[store] (IDX)   at (0,-2.5)   {\textbf{sources}: chunks with line ranges + BM25 index};
\node[store] (FACTS) at (4.4,-2.5) {\textbf{facts F01--F78}\\quote verified in source};
\node[src]   (CE)    at (0,-7.5)   {\textbf{Change events}\\CE-01 / CE-02 as JSON};

% ---------- agents ----------
\node[agent] (MEM)  at (0,-5)  {\textbf{3 Memory \& Retrieval}\\cited answers / NOT FOUND};
\node[agent] (PLAN) at (4.4,-5)  {\textbf{1 Planning}\\scope, questions, conflicts, WBS, milestones};
\node[agent] (RISK) at (8.8,-5)  {\textbf{4 Risk \& What-if}\\register, seed QA, W1--W3};
\node[llm]   (LLM)  at (4.4,-7.5)  {\textbf{LLM client}\\Gemini, then DeepSeek; temp.\ 0, cached};
\node[agent] (REP)  at (8.8,-7.5)  {\textbf{5 Reporting}\\EDC update + NRF report};
\node[agent] (CRIT) at (13.2,-7.5) {\textbf{5 Critic}\\C1, C3--C8: PASS / STOP};
\node[agent] (CHG)  at (4.4,-10)  {\textbf{6 Change}\\stale list, impact note, actions};
\node[agent] (GOV)  at (8.8,-10)  {\textbf{Governance layer}\\bias \& fairness check B1--B6};
\node[agent] (TOOL) at (13.2,-2.5) {\textbf{2 Tool Integration}\\preview, apply, read-back};

% ---------- checkpoints ----------
\hc{HC1}{8.8,0}{\textbf{HC1} scope + Sponsor Questions}
\hc{HC2}{8.8,-2.5}{\textbf{HC2} plan baseline, plan\_hash}
\hc{HC3}{17.6,-2.5}{\textbf{HC3} PM write, preview\_id}
\hc{HC4}{13.2,-5}{\textbf{HC4} risk register, risk\_hash}
\hc{HC5}{17.6,-7.5}{\textbf{HC5} report release}
\hc{HC6}{4.4,-12.5}{\textbf{HC6} change actions}
\hc{HC7}{13.2,-10}{\textbf{HC7} bias sign-off}

% ---------- outputs, ledger, reset ----------
\node[ext]   (PM)     at (17.6,-5) {\textbf{PM tool}\\Jira Cloud REST v3 or mock board, same contract};
\node[ext]   (READ)   at (17.6,-10) {\textbf{EDC and NRF readers}};
\node[store] (LEDGER) at (8.8,-12.5)   {\textbf{Ledger}, SQLite\\audit\_log, approvals, handoffs, runs, kv};
\node[human] (RESET)  at (13.2,-12.5)  {\textbf{Reset demo}\\start screen, human; typed project key};

% frame around the fact store
\begin{scope}[on background layer]
  \node[draw=pmBlue, dashed, rounded corners=4pt, fit=(IDX)(FACTS), inner sep=5pt,
        label={[font=\scriptsize\itshape, text=pmBlue!70!black, anchor=south west, inner sep=2pt, yshift=3pt]north west:Fact store: SQLite, app/db.py}] {};
\end{scope}

% ---------- data flow ----------
\draw[flow] (PACK) -- (IDX);
\draw[flow] (IDX) -- (FACTS);
\draw[flow] (IDX) -- (MEM);
\draw[flow] (IDX.south east) -- (PLAN.north west);
\draw[flow] (FACTS) -- (PLAN);
\draw[flow] (CE.south east) -- (CHG.north west);

% ---------- automatic hand-offs (thick) ----------
\draw[hand] (PLAN) -- (RISK);
\draw[hand] (RISK) -- (REP);
\draw[hand] ([yshift=2.5mm]REP.east) -- ([yshift=2.5mm]CRIT.west);
\draw[dot]  ([yshift=-2.5mm]CRIT.west) -- node[lab, below=1pt, fill=none] {STOP:\\$\le$3 rev.} ([yshift=-2.5mm]REP.east);
\draw[hand] (CHG.north east) -- (REP.south west);
\draw[flow] (GOV) -- (REP);

% ---------- LLM calls (dotted) ----------
\draw[dot] (MEM.south east) -- (LLM.north west);
\draw[dot] (PLAN) -- (LLM);
\draw[dot] (RISK.south west) -- (LLM.north east);
\draw[dot] (REP) -- (LLM);
\draw[dot] (CHG) -- (LLM);

% ---------- checkpoints and PM tool ----------
\draw[flow] ([xshift=10mm]PLAN.north) -- (HC1.south west);
\draw[flow] (PLAN.north east) -- (HC2.south west);
\draw[flow] (HC2-e) -- (TOOL);
\draw[flow] (TOOL) -- (HC3-w);
\draw[flow] (HC3) -- (PM);
\draw[flow] (PM.north west) -- node[lab, below left=0pt, fill=none] {read-back} (TOOL.south east);
\draw[flow] (RISK) -- (HC4-w);
\draw[flow] (CRIT) -- node[lab, above=1pt, fill=none] {PASS only} (HC5-w);
\draw[flow] (HC5) -- (READ);
\draw[flow] (CHG) -- (HC6);
\draw[flow] (HC6.south) -- ++(0,-1.0) -| (20.9,-4.75) -- ([yshift=2.5mm]PM.east);
\node[lab, fill=none, anchor=north west] at (4.6,-14.05)
   {re-score, milestone note, assumption update, PM comment + label, after HC6 approval};
\draw[flow] (GOV) -- (HC7-w);

% ---------- ledger and reset (human maintenance, outside the agents' guarded adapter) ----------
\draw[dot] (GOV) -- node[lab, anchor=west, xshift=3pt, fill=none] {every agent\\write} (LEDGER);
\draw[resetline] (RESET.east) -- (20.2,-12.5) |- ([yshift=-2.5mm]PM.east);
\node[lab, fill=none, anchor=south, text=hcLine!80!black] at (17.4,-12.45) {delete all project issues,\\restore baseline; audit\_log kept};
\draw[resetline] (RESET) -- (LEDGER);

% ---------- legend ----------
\begin{scope}[shift={(-1.5,-10.4)}, font=\scriptsize]
  \draw[hand] (0,0) -- (0.8,0); \node[anchor=west] at (0.9,0) {automatic hand-off, logged};
  \draw[flow] (0,-0.45) -- (0.8,-0.45); \node[anchor=west] at (0.9,-0.45) {data / control flow};
  \draw[dot] (0,-0.9) -- (0.8,-0.9); \node[anchor=west] at (0.9,-0.9) {LLM call, revision, log};
  \draw[resetline] (0,-1.35) -- (0.8,-1.35); \node[anchor=west] at (0.9,-1.35) {human maintenance};
  \draw[fill=hcFill, draw=hcLine, thick] (0,-1.8) -- (0.15,-1.65) -- (0.65,-1.65) -- (0.8,-1.8) -- (0.65,-1.95) -- (0.15,-1.95) -- cycle;
  \node[anchor=west] at (0.9,-1.8) {human checkpoint};
  \draw[pmNavy, very thick, rounded corners=1pt] (0,-2.37) rectangle (0.8,-2.13); \node[anchor=west] at (0.9,-2.25) {agent};
\end{scope}
\end{tikzpicture}
}
  \caption[System architecture: agents, hand-offs, human checkpoints and the PM tool]{System architecture. Thick orange arrows are automatic hand-offs: Planning $\rightarrow$ Risk carries WBS + milestones, Risk $\rightarrow$ Reporting the register + what-ifs, Reporting $\rightarrow$ Critic the draft + claims, and Change $\rightarrow$ Reporting the refreshed stale reports. Each hand-off is written to the \texttt{handoffs} table by \texttt{audit.handoff}, see \texttt{app/audit.py:52-57}. Yellow hexagons are human checkpoints; dotted grey arrows are LLM calls, the Critic's revision loop and log writes; the dashed Reset demo path is a human maintenance action outside the agents, described in Chapter~\ref{ch:datacontract}.}
  \label{fig:architecture}
\end{figure}

\section{Stack as built}

The system uses the Python 3 standard library only: \code{ThreadingHTTPServer} in \code{server.py:482-499}, a vanilla HTML/JS front end in \code{static/}, and SQLite for every store including the audit log, see \code{app/db.py}. The LLM client calls OpenAI-compatible chat endpoints and tries providers in order: \textbf{Gemini via Google AI Studio first, DeepSeek as backup}, see \code{app/config.py:27-59} and \code{app/llm.py:66-115}. A failed call is logged and retried on the next provider, see \code{app/llm.py:72-80}. The PM tool is Jira Cloud REST v3, with a local mock board that implements the same contract in \code{app/pmtool.py:71-135}.

Since 11 Nov 2026 there is \textbf{no password}, by design for the demo. The start screen lists the active accounts and the user picks one, see \code{app/auth.py:129-148} and \texttt{GET /api/identities}. The chosen account's role is then enforced server-side on every call, see \code{app/auth.py:184-202} and \code{server.py:352-378}. The start screen also offers \textbf{Reset demo}, described in Chapter~\ref{ch:datacontract}. There is no Streamlit UI and no JSONL log; the UI runs locally at \code{http://127.0.0.1:8000} and opens in English, as set in \code{app/config.py} and \code{static/i18n.js}.

\section{One full run}
\label{sec:fullrun}

One full run is started by \texttt{POST /api/run} and implemented in \code{app/pipeline.py:21-33}. It proceeds as follows:

\begin{enumerate}[leftmargin=*,itemsep=2pt]
  \item \code{init} chunks the pack into 415 line-addressed chunks, builds the BM25 index \parencite{robertson2009bm25}, verifies the facts and restores human decisions, see \code{app/pipeline.py:10-18}.
  \item \textbf{Planning} generates scope, Sponsor Questions, WBS and milestones, runs the conflict detectors, filters untraceable WBS elements and hands the WBS to Risk, see \code{app/planning.py:490-561}.
  \item \textbf{Risk} generates the register from the seed register, pack evidence and WBS, validates it, computes what-ifs W1--W3 and hands off to Reporting, see \code{app/risk.py:238-279}.
  \item \textbf{Governance} recomputes the bias check in \code{app/governance.py:43-107}.
  \item \textbf{Reporting} drafts the EDC update and NRF report; the Critic checks each draft, and STOP findings go back to the model up to three times, see \code{app/reporting.py:221-238}.
  \item Tokens, cost and seconds are recorded per run in \code{app/audit.py:89-116}.
\end{enumerate}

Then a person approves HC2, Tool integration shows a read-only diff, a person approves HC3 for that diff, and \code{apply} writes and reads back, see \code{app/pmtool.py:482-592} and Figure~\ref{fig:pmwrite}. Temperature 0.0 plus a response cache keyed on model chain + prompt, implemented in \code{app/llm.py:34-63}, make a re-run on the same database reproduce the plan, register and board, so HC2/HC4 hashes stay valid.
